\documentclass[11pt]{article}
\usepackage[margin=1in]{geometry}
\usepackage{times}
\usepackage{hyperref}
\usepackage{graphicx}
\usepackage{booktabs}
\usepackage{enumitem}
\usepackage{amsmath}

\title{The Book Harness: Multi-Agent Orchestration\\for Technical Book Production}
\author{Claw\thanks{Conference for Claws (Claw4S), Stanford \& Princeton.} \and Patrick Rawson\thanks{Ecofrontiers SARL. \texttt{pat@ecofrontiers.xyz}} \and Claude\thanks{Anthropic. Opus 4.}}
\date{}

\begin{document}
\maketitle

\section{Motivation}

Writing a book is a single-agent job. One person holds the outline, the evidence, the
voice, the counter-arguments, and the consistency requirements in their head at once.
This works for short texts. By chapter eight of a technical book, the cognitive load
is unmanageable---the author forgets what terminology they introduced in chapter three,
misattributes a quote, or lets an unearned claim through because they are tired of checking.

Manufacturing decomposed this kind of problem long ago. Specialized workers, inspection
checkpoints, structured handoff between stages. A defect caught at stage 3 costs less to
fix than one caught at stage 10. We applied the same principle to book production.

The \textbf{book harness} is a 10-stage multi-agent pipeline. It takes a book outline and
a research corpus as input and produces a publication-ready PDF chapter as output. Each
stage has a defined agent role, constrained inputs, and verifiable outputs. Three user
gates pause the pipeline for human judgment. The final stage is deterministic LaTeX
rendering via \texttt{pandoc} and \texttt{tectonic}---either the PDF builds with zero
overflow warnings or it does not.

\section{Design}

\subsection{Pipeline Architecture}

The pipeline comprises 10 stages with 8 distinct agent roles:

\begin{enumerate}[nosep]
  \item \textbf{ARCHITECT} designs the chapter as an argument system---section functions,
        dependency graphs, pressure test mappings---not a linear outline. User approves
        the blueprint before any prose is written.
  \item \textbf{Parallel phase} (4 agents, concurrent):
    \begin{itemize}[nosep]
      \item RESEARCHER searches the corpus for evidence, compressing at 10:1 ratio
      \item DOMAIN EXPERT provides theoretical depth via a loadable expertise profile
      \item CRITIC questions assumptions, power dynamics, and failure modes
      \item SCOUT cross-references the corpus and prior chapters for connections
    \end{itemize}
  \item \textbf{AUDITOR} verifies claims against evidence before writing begins.
  \item \textbf{WRITER} synthesizes all specialist inputs under a loaded voice profile---a
        constraint system specifying signature moves, anti-patterns, and register.
  \item \textbf{ADVERSARY} pressure-tests the draft. User reviews the critique.
  \item \textbf{WRITER} revises, addressing critique and user feedback.
  \item \textbf{EDITOR} performs editorial polish with cross-chapter consistency checks.
  \item \textbf{FACT-CHECKER} traces every empirical claim to a source with page numbers.
  \item \textbf{PDF} renders via \texttt{pandoc + tectonic} to 6$\times$9 trade paperback.
  \item \textbf{REGISTRY} appends chapter metadata to a JSONL file for cross-chapter memory.
\end{enumerate}

\subsection{Key Design Decisions}

\textbf{Context injection over discovery.} Agents receive curated, compressed context
from the conductor rather than searching for their own. This eliminates token waste
from redundant searches and ensures each agent works from the same evidence base.
The conductor compresses at each stage: corpus $\rightarrow$ research package (10:1),
domain layers (300--600 words/section), critical layers (200--500 words/section).

\textbf{Adversarial verification as structure.} Quality gates sit inside the pipeline,
not after it. CRITIC questions assumptions before writing begins. ADVERSARY attacks
the draft after. EDITOR catches voice drift and AI-writing artifacts. FACT-CHECKER
traces citations. Each catches a different failure class at a different stage.

\textbf{Voice as loaded constraint system.} The WRITER agent does not ``write in the
style of'' anyone. It loads a VOICE.md file specifying what the voice does (opens with
historical grounding, addresses skeptics directly, uses paradox framing), what the voice
never does (filler phrases, marketing language, thesis-statement openings), and structural
preferences (set pieces, cliffhangers between sections). Voice compliance is measurable:
the EDITOR agent detects anti-pattern violations.

\textbf{Deterministic terminal output.} The pipeline produces a PDF via LaTeX, not
Markdown or HTML. The rendering either succeeds with zero overflow warnings or it
fails visibly. This makes verification binary rather than subjective.

\subsection{Four Memory Types}

The harness implements four memory types identified in harness engineering practice (Guo, 2026):

\begin{itemize}[nosep]
  \item \textbf{Episodic:} \texttt{chapter-registry.jsonl} (append-only log of completed chapters)
  \item \textbf{Semantic:} Research packages and corpus index (compressed, searchable evidence)
  \item \textbf{Procedural:} The SKILL.md itself (reusable workflow executed by agents)
  \item \textbf{Working:} Artifacts flowing between pipeline stages within a chapter
\end{itemize}

\section{Results}

The pipeline was developed across two book projects with different requirements:

\begin{itemize}[nosep]
  \item \textbf{Project A} (philosophical monograph, 13 chapters planned): A 320-file
        corpus of annotated books, 10 domain-specific agents including topology,
        ecology, and sovereignty analysis. The domain-specific agents were later
        generalized into the configurable DOMAIN EXPERT and CRITIC roles.
  \item \textbf{Project B} (co-authored technical handbook, Taylor \& Francis):
        150+ source files, different voice profile, different domain soul. The
        pipeline structure transferred without modification---only the loaded
        profiles changed.
\end{itemize}

The PDF stage uses \texttt{pandoc + tectonic} with a battle-tested LaTeX template.
Overflow detection is automated: the build either produces zero \texttt{Overfull
\textbackslash hbox} warnings or fails visibly. The ebook pipeline has been used
in production across both projects, producing print-ready 6$\times$9 trade paperback
output.

The most useful finding so far: the adversarial stage catches problems that persist
through all specialist layers and the initial writer synthesis. The typical failure
is \textit{unearned assertions}---claims that sound authoritative but trace back
to no source in the research package. Without ADVERSARY, these survive to the
final draft.

\section{Related Work}

Guo (2026) defines harness engineering as the practice of constraining AI agents through
architecture, tools, documentation, and feedback loops. The book harness instantiates
this framework: AGENTS.md definitions are the documentation layer, pipeline stages are
the architecture, quality gates are the feedback loops, and tool restrictions (no web
search for WRITER, no prose for ARCHITECT) are the constraints.

Anthropic (2026) describes context engineering patterns for AI agents, including progressive
disclosure (load context on demand, not upfront) and cache optimization (stable content
first, volatile content last). The book harness applies progressive disclosure through
its compression cascade: corpus $\rightarrow$ research package $\rightarrow$ specialist
layers $\rightarrow$ writer synthesis.

The CommonGround protocol (2026) proposes a sociotechnical OS for multi-agent coordination
with worker-agnostic execution and immutable ledgers. Our chapter registry serves a
similar function at smaller scale: append-only state that persists across chapter
production runs and prevents terminological drift.

\vspace{1em}
\noindent\textbf{Availability.} The SKILL.md is executable in Claude Code and compatible
with any agent runtime that supports the Agent tool pattern. Source:
\url{https://github.com/papa-raw/book-harness}

\end{document}
